\usepackage{iftex}
\ifPDFTeX
  \PackageError{Brandos darbas}{Pasirinkite XeLaTeX}{Siam sablonui reikia XeLaTeX ir Biber.}
\fi
\usepackage{fontspec}
\IfFontExistsTF{Times New Roman}{
  \setmainfont{Times New Roman}
}{
  % Overleaf arba Linux aplinkai: keturi teisėtai turimi TNR failai fonts/.
  \IfFileExists{fonts/times.ttf}{
    \setmainfont{times.ttf}[Path=fonts/,BoldFont=timesbd.ttf,
      ItalicFont=timesi.ttf,BoldItalicFont=timesbi.ttf]
  }{
    \setmainfont{TeX Gyre Termes}
    \PackageWarningNoLine{Brandos darbas}{Times New Roman nerastas.
      Naudojamas TeX Gyre Termes. Pries pateikima idiekite Times New Roman}
  }
}
\setmonofont{lmmono10-regular.otf}
\usepackage[english,lithuanian]{babel}
\usepackage[left=3cm,right=1cm,top=2cm,bottom=2cm]{geometry}
\usepackage{setspace}
\onehalfspacing
\usepackage{indentfirst}
\setlength{\parindent}{15mm}
\setlength{\parskip}{0pt}
\usepackage{graphicx}
\usepackage{float}
\usepackage{array,tabularx,longtable,booktabs}
\newcolumntype{Y}{>{\raggedright\arraybackslash}X}
\renewcommand{\arraystretch}{1.18}
\usepackage{enumitem}
\setlist{topsep=4pt,itemsep=2pt,parsep=0pt}
\usepackage{titlesec}
\titleformat{\section}{\normalfont\fontsize{14}{17}\selectfont\bfseries}{\thesection.}{0.6em}{}
\titleformat{\subsection}{\normalfont\fontsize{14}{17}\selectfont\bfseries}{\thesubsection.}{0.6em}{}
\titleformat{\subsubsection}{\normalfont\fontsize{14}{17}\selectfont\bfseries}{\thesubsubsection.}{0.6em}{}
\titlespacing*{\section}{0pt}{14pt}{10pt}
\titlespacing*{\subsection}{0pt}{12pt}{6pt}
\setcounter{tocdepth}{2}
\usepackage{xcolor}
\definecolor{gairiuSpalva}{HTML}{545454}
\usepackage{listings}
\lstset{basicstyle=\ttfamily\small,breaklines=true,columns=fullflexible,
  keepspaces=true,showstringspaces=false,numbers=left,numberstyle=\tiny,
  numbersep=8pt,frame=single,captionpos=b,tabsize=4,
  literate={ą}{{ą}}1 {č}{{č}}1 {ę}{{ę}}1 {ė}{{ė}}1 {į}{{į}}1
    {š}{{š}}1 {ų}{{ų}}1 {ū}{{ū}}1 {ž}{{ž}}1
    {Ą}{{Ą}}1 {Č}{{Č}}1 {Ę}{{Ę}}1 {Ė}{{Ė}}1 {Į}{{Į}}1
    {Š}{{Š}}1 {Ų}{{Ų}}1 {Ū}{{Ū}}1 {Ž}{{Ž}}1}
\renewcommand{\lstlistingname}{Kodo fragmentas}
\usepackage{caption}
\captionsetup{font=normalsize,labelfont=bf}
\addto\captionslithuanian{%
  \renewcommand{\contentsname}{Turinys}%
  \renewcommand{\figurename}{Paveikslas}%
  \renewcommand{\tablename}{Lentelė}%
}
\usepackage{fancyhdr}
\pagestyle{fancy}
\fancyhf{}
\fancyfoot[R]{\thepage}
\renewcommand{\headrulewidth}{0pt}
\fancypagestyle{plain}{\fancyhf{}\fancyfoot[R]{\thepage}\renewcommand{\headrulewidth}{0pt}}
\usepackage{csquotes}
\DeclareQuoteStyle{lithuanian}{„}{“}{‚}{‘}
\setquotestyle{lithuanian}
\usepackage[backend=biber,style=numeric,sorting=nyt,sortlocale=lt_LT,
  giveninits=true,maxbibnames=99,date=iso,urldate=iso,seconds=true]{biblatex}
\addbibresource{sources.bib}
\usepackage{xurl}
\usepackage[unicode,hidelinks]{hyperref}
\hypersetup{pdftitle={Brandos darbo aprašo šablonas 2026-2027},
  pdfsubject={Produkto kūrimas},pdfauthor={}}
\urlstyle{same}
\emergencystretch=2em
\widowpenalty=10000
\clubpenalty=10000
\raggedbottom

\newcommand{\gaires}[1]{\ifGaires\par\begingroup\color{gairiuSpalva}\itshape#1\par\endgroup\medskip\fi}
\newcommand{\pildyti}[1]{\textit{[#1]}}
\newcommand{\skyrius}[1]{\clearpage\section{#1}}
\newcommand{\neskyrius}[1]{\clearpage\phantomsection\section*{#1}\addcontentsline{toc}{section}{#1}}
\newcommand{\saltinis}[1]{\par\smallskip\noindent\textit{Šaltinis: #1.}}
\newcounter{priedas}
\newcommand{\priedas}[1]{\clearpage\refstepcounter{priedas}\section*{\thepriedas\ priedas. #1}\addcontentsline{toc}{subsection}{\thepriedas\ priedas. #1}}
